% NSF MSPRF Project Description (NSF 23-603) -- DRAFT
%
% Limit: 5 pages. Required content: introduction/background; past accomplishments;
% research objectives, methods, and significance; how the fellowship enhances career
% development; justification of sponsoring scientist and host; and a separate section
% labeled "Broader Impacts". Write for both specialists and non-specialists.
%
% PAPPG formatting: Computer Modern at >= 11pt, margins >= 1in, no URLs in the
% Project Description. References Cited is uploaded as a separate document (no page
% limit); it is compiled after a \clearpage here only for convenience.
%
% Comments marked NOTE flag passages that were rephrased from research_statement.tex
% to avoid the issues discussed earlier; TODO marks missing information.

\documentclass[11pt]{article}

\usepackage[margin=1in]{geometry}
\usepackage{amsmath, amssymb, amsthm}
\usepackage{mathtools}
\usepackage{microtype}
\usepackage{enumitem}
\usepackage[compact]{titlesec}
\usepackage[hidelinks]{hyperref}
\usepackage[nameinlink, capitalize]{cleveref}

\titleformat{\section}{\large\bfseries}{\thesection.}{0.5em}{}
\titleformat{\subsection}{\normalsize\bfseries}{\thesubsection.}{0.5em}{}
\titlespacing*{\section}{0pt}{1.2ex plus .4ex minus .2ex}{0.5ex plus .2ex}
\titlespacing*{\subsection}{0pt}{0.9ex plus .3ex minus .2ex}{0.4ex plus .2ex}
\setlist{nosep, leftmargin=*}

% Block paragraphs: no first-line indent, extra vertical space between paragraphs.
\setlength{\parindent}{0pt}
\setlength{\parskip}{0.3\baselineskip plus 2pt minus 1pt}

\newcommand{\Q}{\mathbb{Q}}
\newcommand{\C}{\mathbb{C}}
\newcommand{\Z}{\mathbb{Z}}
\newcommand{\Gm}{\mathbb{G}_m}
\newcommand{\bbA}{\mathbb{A}}
\newcommand{\bbP}{\mathbb{P}}
\newcommand{\cA}{\mathcal{A}}
\newcommand{\cM}{\mathcal{M}}
\newcommand{\cO}{\mathcal{O}}
\newcommand{\cP}{\mathcal{P}}
\newcommand{\MT}{\mathbf{MT}}
\newcommand{\Sp}{\mathrm{Sp}}
\newcommand{\GSp}{\mathrm{GSp}}
\newcommand{\SL}{\mathrm{SL}}
\newcommand{\IH}{\mathrm{IH}}
\newcommand{\Kum}{\mathrm{Kum}}

\newtheoremstyle{compact}{3pt}{3pt}{\itshape}{}{\bfseries}{.}{.5em}{}
\theoremstyle{compact}
\newtheorem{theorem}{Theorem}
\renewcommand{\thetheorem}{\Alph{theorem}}
\newtheorem{question}{Question}
\newtheorem*{theorem*}{Theorem}
\newtheorem*{conjecture*}{Conjecture}
\crefname{question}{Question}{Questions}

% TODO: working title; final choice pending.
\newcommand{\projecttitle}{The Hodge theory of Lagrangian fibrations, 
with applications to hyper-K\"ahler geometry}

\begin{document}

\begin{center}
  {\Large\bfseries \projecttitle}\\[5pt]
  {\bf Applicant: }Edward A. Varvak, University of Chicago \\ 
  {\bf Sponsoring scientist}: Davesh Maulik, Massachusetts Institute of Technology
\end{center}

%==========================================================================================
\section{Introduction and background}
%==========================================================================================

Compact K\"ahler manifolds with trivial canonical bundle, the higher-dimensional
analogues of elliptic curves and K3 surfaces, are central objects in algebraic geometry
and mathematical physics. By the Beauville--Bogomolov decomposition theorem
\cite{beauville83,bogomolov74}, every such manifold is, up to a finite \'etale cover, a
product of complex tori, Calabi--Yau manifolds, and \emph{hyper-K\"ahler manifolds}:
simply connected compact K\"ahler manifolds $X$ whose holomorphic $2$-forms are spanned by
a single symplectic form~$\sigma$. The decomposition extends to varieties with mild (klt)
singularities \cite{bgl22}, with hyper-K\"ahler manifolds replaced by their singular
analogues, irreducible symplectic varieties.

% NOTE: reframed relative to research_statement.tex (HK vs. CY deformation spaces).
Hyper-K\"ahler manifolds are governed by Hodge theory to a remarkable degree: by the global
Torelli theorem \cite{verbitsky13}, a deformation class is described by an explicit
arithmetic quotient of a period domain, so a single example determines the whole class.
The fundamental open problem is instead to determine which deformation classes
exist. In each dimension $2n$ only two families are known, $K3^{[n]}$ and $\Kum_n$, along
with two sporadic examples, $OG6$ and $OG10$, in dimensions $6$ and $10$. The guiding
conjecture in this direction, the hyper-K\"ahler SYZ conjecture, was recently proven in a
manuscript released by OpenAI \cite{openai-syz26}, and independently by Engel--Mauri
\cite{engel-mauri26}.

\begin{theorem*}[Hyper-K\"ahler SYZ]
Let $L$ be a nontrivial nef line bundle on a hyper-K\"ahler manifold $X$ which is isotropic
for the Beauville--Bogomolov--Fujiki form $q_X$. Then $L$ is semiample, and a power of $L$
defines a \emph{Lagrangian fibration}: a surjective morphism $f\colon X\to B$ with connected
fibers of dimension $n=\frac12\dim X$ on which $\sigma$ vanishes.
\end{theorem*}

When $b_2(X)\ge5$, as for all known examples, $q_X$ has isotropic integral classes, so every
such deformation class contains a Lagrangian fibration. Combined with the boundedness results
of Engel--Filipazzi--Greer--Mauri--Svaldi \cite{engel25}, it follows that there are only
finitely many deformation classes with $b_2\ge5$ in each dimension
\cite{openai-syz26,engel-mauri26}. Lagrangian fibrations are therefore the natural gateway to
classification. They are also highly constrained:
if $f: X\to B$ is a Lagrangian fibration of a hyper-K\"ahler manifold of dimension
$2n$, then its smooth fibers are $n$-dimensional abelian varieties \cite{matsushita99};
the base $B$ is a $\Q$-factorial klt Fano variety of Picard rank one with the intersection
cohomology of $\bbP^n$ \cite{matsushita99,shen-yin22}, and $B\cong\bbP^n$ whenever $B$ is
smooth \cite{hwang08,greb-lehn14,bakker-schnell23} or, by a second recent OpenAI manuscript,
whenever $X$ is projective \cite{openai-bases26}, proving a conjecture of Matsushita in the
projective case; and $f$ is either isotrivial or of \emph{maximal variation}
\cite{bakker25}, meaning that the period map $\cP\colon B^\circ\to\cA_n$, which sends a
point of the smooth locus $B^\circ\subset B$ to the isomorphism class of the fiber $X_b$ in
the moduli space $\cA_n$ of polarized abelian varieties, has $n$-dimensional image.

The goal of this proposal is to determine where the period map lands, i.e. which
abelian varieties, and which families of them, occur as fibers of Lagrangian
fibrations. We also hope to use this Hodge-theoretic information to attack structural questions:
the classification of Lagrangian fibrations whose fibers decompose
(\Cref{q:classification,q:k3n}), the extension of the theory to non-compact settings such as
Hitchin systems (\Cref{q:noncompact,q:hitchin}), and the construction of rational curves and
uniruled divisors on hyper-K\"ahler manifolds (\Cref{HKBM}).

%==========================================================================================
\section{Past results}
%==========================================================================================

My thesis work, supervised by Benjamin Bakker, develops the Hodge theory of Lagrangian
fibrations. Let $X$ be a \emph{primitive symplectic variety} of dimension $2n$: a compact
symplectic variety with $H^1(X,\cO_X)=0$ and a unique symplectic form up to scaling, such as
a hyper-K\"ahler manifold. Let $f: X\to B$ be a Lagrangian fibration, whose smooth
fibers are abelian varieties \cite{schwald20}. Let $f^\circ: X^\circ\to B^\circ$ be the
smooth part of $f$ and $V=R^1f^\circ_*\Z$ the associated variation of Hodge structures,
viewed as a representation of $\pi_1(B^\circ)$.

\begin{theorem}[\cite{varvak25}]\label{thm:monodromy}
If $f$ is not isotrivial, then $V_\C$ is irreducible. If $f$ is isotrivial, then the smooth
fibers are isogenous to $E^n$ for an elliptic curve $E$ (first proven in
\cite{kim-laza-martin24}), and $V_K=U_1\oplus U_2$ splits into two geometrically irreducible
summands, where $K$ is the endomorphism field of $E$.
\end{theorem}

% NOTE: "over Q" added.
\begin{theorem}[\cite{varvak26}]\label{thm:mt}
If $f$ is not isotrivial, then for very general $b\in B$ the Mumford--Tate group of $X_b$ is
isogenous over $\Q$ to $\Gm\cdot(\Sp_{2h})^{n/h}$ for some $h\mid n$. Equivalently, $X_b$ is
isogenous to a product $A_1\times\cdots\times A_{n/h}$ of pairwise non-isogenous
$h$-dimensional abelian varieties, each \emph{Hodge-generic}, i.e.\ carrying no Hodge
classes beyond those forced by its polarization.
\end{theorem}

\begin{theorem}[\cite{varvak26}]\label{thm:general}
Let $\mathcal{X}$ be a deformation class of hyper-K\"ahler manifolds with $b_2\ge7$ which
contains a Lagrangian fibration. If $X\in\mathcal{X}$ is general among members admitting a
Lagrangian fibration $f\colon X\to B$, then the very general fiber of $f$ is Hodge-generic:
$\MT(X_b)=\GSp_{2n}$.
\end{theorem}

These results sharpen Bakker's theorem: the period image is not only $n$-dimensional but
Hodge-generic either in $\cA_n$ or in a product locus $\cA_h^{n/h}$, and in the general case
it is the former. The proof of \Cref{thm:mt} hinges on the interplay between the connected
monodromy group, which preserves each factor $A_i$, and the finite part of the monodromy,
which by \Cref{thm:monodromy} must permute the factors transitively. The hypothesis
$b_2\ge7$ holds for every known deformation type, and by the SYZ theorem every such
deformation class contains a Lagrangian fibration.

\subsection*{Personal accomplishments}

I've been invited to speak about my research at the Institute for Advanced Study's
Research Seminar (November 2025), GTA Philadelphia (May 2026), and the upcoming 
Special Session on Higgs Bundles at JMM 2027. I've received the Kuga--Sah Award for
Excellence in Mathematics from the Stony Brook University Department of Mathematics (2022). 
I have received the UIC RTG grant for six semesters throughout my time at UIC. 
I received the Bernard Kurtin Award for research in mathematics from the UIC Department of 
Mathematics, Statistics, and Computer Science (2025). 

%==========================================================================================
\section{Research objectives, methods, and significance}
%==========================================================================================

We will break our research plan into three primary objectives, which we order by the
expected difficulty of the task in question.
\begin{enumerate}[label=\textbf{RO \arabic*.}]
  \item Determine decompositions of the Mumford--Tate group occur for Lagrangian fibers.
  \item Understand the Hodge theory of non-compact Lagrangian fibrations.
  \item Work towards the Bogomolov--Mumford conjecture for hyper-K\"ahler manifolds.
\end{enumerate}

In this section, we will go through each of them, laying out a motivation and an 
approach for how to address them.

\subsection{Which decompositions of the Mumford--Tate group occur?}

\Cref{thm:mt} leaves open which values of $h$ actually occur. If $S\to\bbP^1$ is a
non-isotrivial elliptic K3 surface, the induced map $S^{[n]}\to(\bbP^1)^{(n)}\cong\bbP^n$ is
a Lagrangian fibration whose fibers are isogenous to products of $n$ elliptic curves, so
$h=1$. If singularities are allowed, every $h$ occurs: given a Lagrangian fibration
$Y\to\bbP^h$ with Hodge-generic fibers (\Cref{thm:general}), the symmetric power
$Y^{(k)}\to(\bbP^h)^{(k)}$ is a Lagrangian fibration of a primitive symplectic variety with
$\MT(X_b)\sim\Gm\cdot(\Sp_{2h})^k$. For $h\ge2$, however, the base $(\bbP^h)^{(k)}$ does not
have the intersection cohomology of projective space (e.g.\
$\dim\IH^4((\bbP^2)^{(2)})=2$), so it cannot be the base of a Lagrangian fibration of a
hyper-K\"ahler manifold.

\begin{question}\label{q:classification}
Is there a hyper-K\"ahler manifold (or irreducible symplectic variety) with a Lagrangian
fibration such that $\MT(X_b)\sim\Gm\cdot(\Sp_{2h})^{n/h}$ with $1<h<n$?
\end{question}

\begin{question}\label{q:k3n}
If $X$ is hyper-K\"ahler and the very general fiber is not Hodge-generic, i.e.\ $h<n$, is
$X$ of $K3^{[n]}$-type?
\end{question}

I expect a negative answer to \Cref{q:classification}: whenever the fibers are not
Hodge-generic, the fibration should be governed by families of elliptic curves, as it is
for $S^{[n]}$. A complete answer would show that decomposable Lagrangian fibrations
are exactly those built from elliptic fibrations and give a Hodge-theoretic characterization
of the $K3^{[n]}$ deformation type.


% NOTE: this approach paragraph was added in the draft; check it against your intended
% strategy.
\emph{Approach.} Suppose $h<n$, set $k=n/h$, and pass to the finite \'etale cover
$B'\to B^\circ$ over which $V$ splits as $V_1\oplus\cdots\oplus V_k$. The symplectic form
identifies $T_bB$ with $H^{1,0}(X_b)$, and by Donagi--Markman \cite{donagi-markman96} the
differential of the period map is a symmetric cubic form on $T_bB$. Because the period map
respects the splitting, this cubic form is a sum of cubic forms on the summands
$E_i\subset T_bB$ corresponding to the $H^{1,0}(V_{i,b})$. Consequently each partial period
map $\cP_i\colon B'\to\cA_h$ depends only on the $E_i$-directions, the $E_i$ integrate to
algebraic foliations, and $(\cP_1,\dots,\cP_k)$ maps $B'$ dominantly and generically finitely
onto a product $Z_1\times\cdots\times Z_k$ of $h$-dimensional varieties. This is exactly the
structure of $(\bbP^1)^n\to\bbP^n$ in the example above. I plan to extend this structure
across the discriminant locus and combine it with the constraints on $B$ (Fano, Picard rank
one, $\IH^*(B)\cong H^*(\bbP^n)$, and now $B\cong\bbP^n$ when $X$ is a projective
hyper-K\"ahler manifold \cite{openai-bases26}) and an analysis of the branching of $B'\to B$
to rule out $1<h<n$. When $h=1$, I will use the factors to construct an elliptic K3 surface $S$ and
compare $X$ with a moduli space of sheaves on $S$, identifying its deformation type. Recent
work of Sacc\`a on compactifications of Lagrangian fibrations \cite{sacca25} will be an important
input.


\subsection{Lagrangian fibrations beyond the compact case}

Many important Lagrangian fibrations are not compact: integrable systems from mathematical
physics, local models near singular fibers, and above all the Hitchin fibration
$h\colon\cM_{\mathrm{Dol}}\to\bbA$ on the moduli space of Higgs bundles on a curve $C$
\cite{hitchin87}, which is central to non-abelian Hodge theory and the geometric Langlands
program. My compact methods use Bakker's dichotomy, which fails here: $h$ is equivariant for
the scaling action $(E,\theta)\mapsto(E,\lambda\theta)$, so its period map is constant along
$\C^*$-orbits and never has maximal variation.

% NOTE: divisibility condition on the h_i removed.
\begin{question}\label{q:noncompact}
Let $f\colon M\to B$ be a proper Lagrangian fibration of a (not necessarily compact)
holomorphic symplectic manifold, of maximal variation. Is the Mumford--Tate group of its
very general fiber isogenous to $\Gm\cdot(\Sp_{2h_1}\times\cdots\times\Sp_{2h_k})$ with
$h_1+\cdots+h_k=n$?
\end{question}

% NOTE: stated for SL_r (for GL_r, R^1h_* has a constant summand H^1(C)).
I expect the foliation and monodromy techniques above to extend to this setting, possibly
under mild hypotheses on the singular fibers. For the Hitchin fibration of
$\SL_r$-Higgs bundles, whose regular fibers are Prym varieties of spectral covers
$\Sigma_a\to C$, I can adapt the methods of \cite{varvak25} to show that $R^1h_*\C$ is
irreducible over the regular locus $\bbA^\circ$, despite the failure of maximal variation.

\begin{question}\label{q:hitchin}
Is the very general $\SL_r$-Hitchin fiber Hodge-generic? If not, which products of
symplectic groups occur as its Mumford--Tate group?
\end{question}

Mumford--Tate groups govern the Hodge classes, and conjecturally the
algebraic cycles, on the fibers. This complements recent work of Maulik--Shen--Yin on
algebraic cycles and Fourier transforms for Hitchin systems \cite{msy-cycles,msy-fourier},
and gives a new Hodge-theoretic invariant of a fundamental integrable system.

\emph{Approach.} Spectral curves form an open subset of a linear system on the
compactified cotangent bundle $\bbP(K_C\oplus\cO_C)$, so Lefschetz theory is available: the
monodromy is generated by Picard--Lefschetz transformations in vanishing cycles, and
Zariski density results $\cite[\mathrm{Lemma\ 4.4.2}^a]{deligne74}$ apply. Since the connected monodromy group is a
normal subgroup of the derived Mumford--Tate group \cite{andre92}, showing that the
monodromy on the Prym part is Zariski-dense in $\Sp_{2d}$, $d=\dim\bbA$, would prove
Hodge-genericity. The difficulties
are that the Hitchin base is only part of a linear system and that the relevant vanishing
cohomology is that of the Prym; here we will try to apply the work of Maulik and his
collaborators on the topology of Hitchin fibrations \cite{dCMS22,maulik-shen24}.


% Commented out: a fourth objective on abelian fibrations of Calabi--Yau varieties. To
% restore it, also list it in the goals paragraph of Section 1, change "three" to "four" at
% the start of Section 3 and add it to the list of objectives there, and mention
% Calabi--Yau fibrations in "Broader scientific impact".
%
% \subsection{Abelian fibrations of Calabi--Yau varieties}\label{sec:cy}
%
% Strict Calabi--Yau varieties $Y$ (simply connected, with $\omega_Y\cong\cO_Y$ and
% $h^{p,0}(Y)=0$ for $0<p<\dim Y$), the other simply connected factors in the
% Beauville--Bogomolov decomposition, often carry fibrations $f\colon Y\to B$ whose general
% fibers are $g$-dimensional abelian varieties. Elliptic Calabi--Yau threefolds, which are
% central to F-theory, and Calabi--Yau threefolds fibered in abelian surfaces
% \cite{schoen88,gross-popescu01} are basic examples, and Calabi--Yau varieties of fixed
% dimension with such fibrations are birationally bounded \cite{engel25}. Part of the
% hyper-K\"ahler picture survives: there is a period map $\cP\colon B^\circ\to\cA_g$, and if
% $f$ has maximal variation ($\cP$ generically finite), the canonical bundle formula
% \cite{ambro05} shows that $B$ is of Fano type. But the symplectic form is gone: $\dim B$ is
% no longer tied to $g$, the period map need not satisfy the Donagi--Markman cubic condition,
% and the restriction map $H^2(Y,\Q)\to H^2(Y_b,\Q)$, which has rank one for Lagrangian
% fibrations \cite{shen-yin22}, now has image $\mathrm{NS}(Y_\eta)_\Q$, where $Y_\eta$ is the
% generic fiber, and this can be large. Accordingly, the naive analogue of \Cref{thm:mt} fails:
% a small resolution of the self-fiber product $S\times_{\bbP^1}S$ of a non-isotrivial rational
% elliptic surface $S\to\bbP^1$ is a strict Calabi--Yau threefold \cite{schoen88} fibered in
% abelian surfaces $E_t\times E_t$, whose period image is not Hodge-generic in any product
% locus.
%
% \begin{question}\label{q:cy}
% Let $f\colon Y\to B$ be an abelian fibration of a strict Calabi--Yau variety with maximal
% variation, and suppose that, as for Lagrangian fibrations, the restriction map
% $H^2(Y,\Q)\to H^2(Y_b,\Q)$ has rank one. Is the image of the period map Hodge-generic in a
% product locus of $\cA_g$? If not, which Mumford--Tate groups occur for the very general
% fiber?
% \end{question}
%
% The rank-one hypothesis holds, for instance, when $\rho(Y)=\rho(B)+1$, and it rules out the
% example above. Even so, I expect a negative answer: it does not rule out fibers with real
% multiplication or with isogenous factors permuted by the monodromy, and nothing replaces the
% Donagi--Markman condition. An answer would show which features of
% Theorems~\ref{thm:monodromy}--\ref{thm:general} come from K-triviality and which require the
% symplectic form, and would constrain the fibers of abelian fibrations of Calabi--Yau
% varieties.
%
% \emph{Approach.} The first test case is a Calabi--Yau threefold $f\colon Y\to\bbP^1$ fibered
% in abelian surfaces. Its period image is a curve in $\cA_2$, and the smallest special
% subvariety containing it is $\cA_2$, a product locus, a Hilbert modular surface, or a special
% curve. Since $\omega_Y\cong\cO_Y$, we have $f_*\omega_{Y/\bbP^1}\cong\cO_{\bbP^1}(2)$, so if
% the local monodromies are unipotent, the Hodge bundle has degree $2$; the Arakelov inequality
% then forces at least four singular fibers, with equality only if the period image is a
% Shimura curve \cite{viehweg-zuo04}. I will combine such degree constraints with the geometry
% of curves on Hilbert modular surfaces to decide which cases occur, using the constructions of
% \cite{schoen88,gross-popescu01} as test cases. In general, I will adapt the proofs of
% \Cref{thm:monodromy,thm:mt} to $V=R^1f^\circ_*\Z$, with the rank-one hypothesis in place of
% the symplectic form: it should again force the finite part of the monodromy to permute
% transitively the isotypic summands of $V$ over a finite \'etale cover of $B^\circ$, reducing
% the problem to the connected monodromy group of each summand, a normal subgroup of its derived
% Mumford--Tate group \cite{andre92}. The new difficulty is that the summands may carry real or
% complex multiplication, which must be ruled out or realized using the global geometry of $Y$.


\subsection{Rational curves and the Bogomolov--Mumford conjecture}\label{HKBM}

By a theorem of Bogomolov and Mumford \cite{mori-mukai83}, every ample linear system on a K3
surface contains a divisor supported on rational curves. The analogue for hyper-K\"ahler
manifolds is closely tied to the study of algebraic cycles on hyper-K\"ahler varieties
\cite{cmp19}; it is known for the deformation types $K3^{[n]}$ and $\Kum_n$ but open in
general.
% TODO: check citations for the K3^[n] and Kum_n cases.

\begin{conjecture*}[Hyper-K\"ahler Bogomolov--Mumford] 
If $X$ is a projective hyper-K\"ahler manifold and $L$ is an ample line bundle on $X$, then
$|L|$ contains a uniruled divisor.
\end{conjecture*}

% NOTE: Ax--Schanuel step rephrased via typical intersections; decomposition rephrased as
% horizontal + vertical.
\emph{Approach.} Let $X$ be general among Lagrangian-fibered members of its deformation
class, which exist by the SYZ theorem, so that by \Cref{thm:general} the period image $\cP(B^\circ)\subset\cA_n$ is
$n$-dimensional and Hodge-generic. Abelian varieties isogenous to a product $E\times A'$
form the Hecke translates of $\cA_1\times\cA_{n-1}$, which has codimension $n-1$, so their
expected intersection with the period image is one-dimensional. Results on the existence
and density of such typical intersections \cite{khelifa-urbanik23,eterovic-scanlon25},
complementing the Ax--Schanuel theorem \cite{bakker-tsimerman19}, which controls the
atypical ones, should produce a dense set of curves $C^\circ\subset B^\circ$ over which the fibers are isogenous to $E_t\times A'_t$.
Over the boundary points of such a curve I expect the elliptic factor to degenerate to a
nodal rational curve while $A'_t$ remains an abelian variety, producing uniruled
subvarieties in the singular fibers of $f$. We need to show that these rational
curves deform to a general polarized deformation of $X$, which no longer carries a
Lagrangian fibration, using the deformation theory of rational curves on hyper-K\"ahler
manifolds \cite{cmp19}. Then we hope to follow the K3 picture, where an ample class is 
a combination $a\sigma+bF$ of a section and a fiber. We will try to decompose a divisor 
in $|L|$ into uniruled horizontal and vertical parts that deform with $X$.


% TODO: I'm not sure that this section should exist, given the absurd pace of advancement 
% in math at present
% \subsection{Timeline}
% \emph{Year 1:} \Cref{q:classification} and \Cref{q:hitchin}. \emph{Year 2:} \Cref{q:k3n} and
% \Cref{q:noncompact}. \emph{Year 3:} \Cref{HKBM}. \Cref{q:classification,q:k3n} build directly
% on my thesis and should produce results early; \Cref{q:noncompact,q:hitchin} and \Cref{HKBM}
% carry more risk and benefit most from the expertise available at MIT.

%==========================================================================================
\section{Career development and choice of sponsoring scientist and host institution}
%==========================================================================================

My goal is a faculty career at a research university, leading a research program on the
Hodge theory of Lagrangian fibrations and its interactions with moduli theory, enumerative
geometry, and non-abelian Hodge theory. My training so far is centered on Hodge theory and
period maps. I have chosen to pursue the  Research Fellowship because full-time research
support suits the long-term, higher-risk parts of this program (such as \Cref{HKBM}) and the
substantial new background they may require. The fellowship will allow me to:
\begin{itemize}
\item learn the sheaf-theoretic and enumerative tools that the proposed questions require: 
perverse filtrations, Fourier transforms on abelian fibrations, moduli of sheaves on K3 surfaces, 
and Gromov--Witten theory; 
\item establish my independence, pursuing the more challenging \Cref{q:hitchin} and 
\Cref{HKBM}, starting collaborations at MIT and in the Boston area, and traveling to 
conferences and research visits supported by the fellowship allowance;
\item develop as a mentor and expositor through the activities described in \Cref{sec:bi}.
\end{itemize}

My sponsoring scientist, Davesh Maulik, is a leading expert on abelian fibrations, 
hyper-K\"ahler varieties, and Hitchin systems, which serve as the backbone of this proposal. 
His work bears directly on each objective. With Shen and Yin he developed the interaction of
Fourier--Mukai transforms with perverse filtrations for abelian fibrations
\cite{msy-fourier}, and with Shen, Yin, and Zhang he proved the D-equivalence conjecture for
hyper-K\"ahler varieties of $K3^{[n]}$-type \cite{msyz24}. These tools, together with his
expertise on moduli of sheaves on K3 surfaces, will help to recognize the
hyper-K\"ahler manifolds in \Cref{q:k3n} as belonging to the deformation space 
\(K3^{[n]}\). 

Professor Maulik is a central figure in the topology
of Hitchin fibrations, including the proof of the $P=W$ conjecture
\cite{dCMS22,maulik-shen24} and the study of algebraic cycles on Hitchin systems
\cite{msy-cycles}. His insight will be invaluable to helping attack \Cref{q:hitchin}. 
Furthermore, his work on Noether--Lefschetz theory
\cite{maulik-pandharipande13,kmps10} and on curves on K3 surfaces \cite{mpt10} sits 
at the interface of Hodge loci and rational curves that \Cref{HKBM} explores. 

% TODO: other MIT/Boston colleagues may be named here later.
MIT offers a large and active community in algebraic geometry and neighboring fields.
Moving from Chicago, where I trained in Hodge theory with Bakker, to MIT, with its strengths
in moduli theory, enumerative geometry, and geometric representation theory, is precisely
the broadening my research program needs.

%==========================================================================================
\section{Broader Impacts}\label{sec:bi}
%==========================================================================================

% \emph{Mentoring young researchers.} At MIT I will mentor high school students through
% PRIMES, which pairs Boston-area students with MIT graduate students and postdocs on research
% projects, and undergraduates through UROP. Classical algebraic geometry offers research
% problems that need little background and lend themselves to computer experimentation; my
% undergraduate thesis, on recovering plane nodal curves from their bitangents, is an example.
% Suitable projects include computing the monodromy action on the 28 bitangents of a varying
% plane quartic and on the vanishing cycles of families of plane curves, and classifying the
% singular fibers of explicit elliptic K3 surfaces, the $n=1$ case of the Lagrangian
% fibrations studied here. Code and results from these projects will be made publicly
% available.

% \emph{Exposition and open resources.} During the fellowship I will expand them into a survey of
% the Hodge theory of Lagrangian fibrations aimed at graduate students---covering monodromy,
% Mumford--Tate groups, and period maps of abelian fibrations---so that students without a
% local expert can enter this rapidly developing area.

\emph{Building research communities.} As a graduate student I co-organized the Graduate
Algebraic Geometry Seminar for two semesters, organized a reading group on the Weil
conjectures, and organized a seminar on Lie groups. At MIT I will organize a learning
seminar on Lagrangian fibrations and Hitchin systems, open to graduate students and postdocs
across the Boston area, and seek support, including from the fellowship allowance, for a
small workshop for early-career researchers. Both will be open to everyone and designed to
reach researchers who lack local access to this area, including those from groups
underrepresented in mathematics. Seminar notes and recorded talks will be posted publicly.
% The workshop will be advertised through broad channels, including the Association for Women
% in Mathematics, the National Association of Mathematicians, and SACNAS; it will open with
% introductory lectures for beginning graduate students; and it will prioritize travel support
% for participants from institutions without a research group in the area, such as primarily
% undergraduate and minority-serving institutions.

\emph{Teaching and outreach.} I am currently teaching Calculus~II at the University of
Chicago as a graduate student lecturer, after five semesters as a teaching assistant at UIC,
a Hispanic-Serving Institution where many students are the first in their families to attend
college. I have also judged the QED competition of the Math Circles of Chicago; in Boston I
will continue this outreach by volunteering with local math circles and competitions such as
HMMT, favoring programs that serve schools with limited access to advanced mathematics. I
will also seek to mentor undergraduates through the MIT Summer Research Program, which brings
students from institutions across the country to MIT to prepare them for graduate study.

\emph{Broader scientific impact.} Lagrangian fibrations and Hitchin systems are algebraic
integrable systems that appear throughout mirror symmetry and mathematical physics; results
on their monodromy and periods will be useful well beyond algebraic geometry.

%==========================================================================================
% References Cited: uploaded separately (no page limit).
%==========================================================================================
\clearpage
% \renewcommand{\refname}{References Cited}
\bibliographystyle{alphaurl}
\bibliography{refs}

\end{document}
